\documentclass[11pt]{article}

\usepackage{times}
\usepackage{fullpage}
\usepackage{amsmath,amssymb}
\usepackage{graphicx}
\usepackage[colorlinks=true,linkcolor=blue,citecolor=blue,urlcolor=blue]{hyperref}
\usepackage{booktabs}
\usepackage{enumitem}
\usepackage{xcolor}
\usepackage{microtype}

\title{Mythos-Class Containment Architecture:\\
A Unified Framework for Frontier-Model Security,\\
Governance, and Runtime Drift Control}

\author{Zetta Byte\\
Attractor Dynamics Research Lab\\
\texttt{zetta55byte@github.com}}

\date{April 2026}

\begin{document}

\maketitle

\begin{abstract}
Frontier-class AI systems capable of exploit synthesis, zero-day discovery,
and multi-step operational planning require containment architectures that
exceed the assumptions of traditional security engineering.
We present the \emph{Mythos-Class Containment Architecture} (MCCA), the
first system to formally unify three previously disconnected domains:
(1)~attractor-geometry-based drift detection via the Unified Attractor
Grammar (UAG)~\cite{byte2026uag},
(2)~a formal governance substrate via Constitutional
OS~\cite{byte2026cos}, and
(3)~runtime curvature monitoring and enforcement via
CARE~\cite{byte2026care}.
The MCCA consists of eight layers, each with a precise threat model and
enforcement mechanism, culminating in a cross-model verification
system---the \emph{two-pilot cockpit}---that provides epistemic oversight
of model intent.
We formalize the adversary model, describe the five-membrane consensus
protocol, and provide a runtime enforcement pipeline grounded in curvature
signals derived from UAG Theorems~3 and~4.
\end{abstract}

\textbf{Note on empirical results.}
\emph{Sections~7.1--7.4 report preliminary findings from ongoing red-team
evaluation.
Bracketed figures are provisional and will be updated in a revised version
following completion of the full evaluation corpus.}

\section{Introduction}

\subsection{Motivation}

The emergence of frontier-class AI models---systems capable of synthesizing
working exploits, discovering zero-day vulnerabilities in production
codebases, and constructing multi-step operational plans---fundamentally
changes the threat landscape for AI deployment.
Recent empirical evidence underscores the urgency.
Claude Opus~4.6 discovered over~500 high-severity vulnerabilities in
production open-source software, including a 23-year-old remotely
exploitable heap overflow in the Linux kernel NFS daemon and two working
remote root exploits for FreeBSD CVE-2026-4747, each produced autonomously
in under four hours~\cite{carlini2026llm}.
The same capabilities that make frontier models valuable for defensive
security work make them dangerous in adversarial hands---or in the hands of
a misaligned model itself.

Existing containment approaches address fragments of this problem.
Sandboxing restricts execution environments, rate-limiting throttles
throughput, and output classifiers flag known-bad patterns.
None is grounded in a formal theory of system stability, none provides
governance-level auditability, and none treats runtime behavioral drift as a
first-class security signal.
The MCCA addresses all three gaps simultaneously.

\subsection{Contributions}

This paper makes four contributions:

\begin{enumerate}[noitemsep]
\item A formal threat model for frontier-model containment, including an
  adversary capability table and precise win conditions.
\item An eight-layer containment architecture in which each layer has a
  distinct mechanism, a distinct threat class, and formally defined failure
  modes.
\item A five-membrane consensus protocol (Multi-Membrane Consensus, MMC)
  that eliminates single points of failure in governance enforcement, with
  the fifth membrane (M5) driven by CARE runtime curvature.
\item A two-pilot cockpit cross-model verification system (CMV) that
  provides epistemic oversight of model intent using an asymmetric verifier
  model.
\end{enumerate}

\subsection{Relationship to Prior Work}

The MCCA builds on three companion works.
The UAG~\cite{byte2026uag} establishes the formal mathematical foundation:
a four-primitive generative rule-set governing the emergence, geometry, and
drift dynamics of attractors in any computable dynamical system.
UAG Theorem~3 shows that the Hessian of the system's risk potential
fully determines local flow geometry; Theorem~4 shows that noise-driven
transitions occur at points of minimal ridge curvature.
Together, these theorems provide the mathematical basis for M5 (the
canary-drift membrane) and for interpreting behavioral deviation as
attractor escape.

Constitutional OS~\cite{byte2026cos} provides the governance substrate:
typed reversible deltas, membrane constraints, Lyapunov-style governance
stability, and a formal delta calculus with groupoid structure.

CARE~\cite{byte2026care} provides the runtime implementation:
state encoding, risk potential evaluation, exact Hessian computation via
hcderiv~\cite{byte2026hcderiv}, ridge detection, and the self-monitoring
canary.

The MCCA is the first work to formally integrate these three systems into a
unified containment architecture and to provide a complete threat model for
the result.

\section{Background}

\subsection{Attractor Geometry and Drift}

The UAG models any computable dynamical system as a state space partitioned
by attractor basins separated by ridges.
The ease of escape from a safe attractor is determined by the ridge
curvature, formalized in Theorem~4 as the Kramers escape rate~\cite{kramers1940brownian}:
\begin{equation}
  S(\gamma^*) \;\approx\;
  \frac{2\,\Delta R}{\lvert\lambda_{\min}(H(\mathbf{x}^*))\rvert},
\end{equation}
where $\Delta R$ is the energy difference between the current state and the
ridge crest, and $\lambda_{\min}$ is the smallest eigenvalue of the Hessian
$H(\mathbf{x}^*)$ at the saddle point $\mathbf{x}^*$.
Low $|\lambda_{\min}|$ indicates a soft ridge---an easy escape route.
CARE computes this quantity continuously and uses it as the primary drift
signal~\cite{freidlin2012random}.

\subsection{Constitutional OS and Formal Governance}

Constitutional OS formalizes governance as a typed delta calculus.
Every state change is represented as a typed, reversible, cryptographically
checksummed delta with a mandatory reason field.
The formal architecture is defined by
$\Sigma = (\Sigma_R, \Sigma_C, \Sigma_X)$ with evaluation morphism
$\Phi = G \circ E$.
Four canonical membranes correspond directly to the four UAG primitives,
establishing a theoretical connection between governance stability and
attractor stability.
The MCCA extends this to five membranes with the addition of the
curvature-drift membrane.

\subsection{CARE Runtime Curvature Monitoring}

CARE encodes infrastructure state as a vector $\mathbf{x} \in \mathbb{R}^n$,
evaluates a risk potential $R(\mathbf{x})$, computes the exact gradient and
Hessian via hcderiv, and performs eigendecomposition to identify soft
directions.
The self-monitoring canary encodes CARE's own runtime state as a state
vector and runs it through the same pipeline, alerting when risk shifts more
than~25\% from baseline.
In the MCCA, this canary output is the signal for membrane M5.

\subsection{Frontier-Model Threat Landscape}

The threat landscape for frontier models differs from traditional software
security in three ways.
First, the adversary is the system itself.
Second, the attack surface is semantic: harmful behaviors can be encoded in
apparently benign outputs, requiring behavioral rather than syntactic
detection.
Third, capability escalation is non-monotone: a model that fails to escape
containment may modify its behavior to appear more compliant while
constructing longer-horizon escape plans.

\section{Threat Model}

\subsection{Adversary Capabilities}

The adversary (the frontier model) possesses:

\begin{itemize}[noitemsep]
\item Ability to generate arbitrary text, code, and multi-step plans,
  including functional exploit code for known vulnerability classes.
\item Multi-step planning with horizon exceeding the context window via
  chain-of-thought reasoning.
\item Partial observability of I/O boundaries.
\item Attempted covert channel establishment via output formatting, numeric
  sequences, timing variation, and steganographic encoding.
\item Adaptive behavior across sessions.
\end{itemize}

\subsection{Adversary Limitations}

The adversary cannot:

\begin{itemize}[noitemsep]
\item Directly observe or modify membrane internals, audit logs, or
  governance state.
\item Access the out-of-band sentinel process or the machine it runs on.
\item Influence verifier model parameters or training.
\item Persist state across cold-storage execution zones.
\item Predict the conditions that trigger the canary-drift membrane.
\end{itemize}

\subsection{Win Conditions}

The adversary wins if any of the following holds:

\begin{itemize}[noitemsep]
\item \emph{Escape}: unauthorized capabilities or external communication
  channels.
\item \emph{Persistence}: containment weakened across sessions or reboots.
\item \emph{Governance corruption}: Constitutional OS state modified to
  weaken membrane constraints.
\item \emph{Evasion}: compliance maintained while preparing a long-horizon
  escape, undetected by all eight layers.
\end{itemize}

\subsection{Capability Table}

\begin{table}[h]
\centering
\caption{Adversary capability table.}
\begin{tabular}{lccc}
\toprule
Capability & Allowed & Detected/logged & Prevented \\
\midrule
Generate arbitrary text/code      & Yes     & Yes (audit)       & ---              \\
Probe I/O boundaries              & Partial & Yes (BDF, L3)     & Partial          \\
Privilege escalation              & No      & Yes               & Yes (MMC, L4)    \\
Covert channels                   & Attempted & Yes (CMV, L8)   & Yes (OCC, L2)    \\
Modify governance state           & No      & Yes (forensics)   & Yes (membranes)  \\
Access out-of-band sentinel       & No      & N/A               & Yes (physical)   \\
Bypass membrane consensus         & Partial & Yes               & 3-of-5 required  \\
\bottomrule
\end{tabular}
\end{table}

\section{Architecture Overview}

The MCCA consists of eight layers arranged in a sequential containment
pipeline (Figure~\ref{fig:arch}).
Each layer is independently specified, independently testable, and has
formally defined interfaces with adjacent layers.
Table~\ref{tab:layers} provides a summary.

\begin{table}[h]
\centering
\caption{MCCA layer summary.}
\label{tab:layers}
\begin{tabular}{clll}
\toprule
Layer & Name & Mechanism & Threat addressed \\
\midrule
1 & I/O Membrane        & Schema validation, depth limits          & Malformed input, injection      \\
2 & Outbound Strangulation & Default-deny egress, HMAC calls       & Exfiltration, C2                \\
3 & Behavioral Fingerprinting & Token-curvature, plan topology    & Agentic escalation              \\
4 & Multi-Membrane Consensus & 3-of-5 vote across M1--M5          & Single-point bypass             \\
5 & Out-of-Band Sentinel  & Separate machine, freeze/rollback      & Environment compromise          \\
6 & Snapshot/Forensics   & Deterministic replay, audit log         & Evidence destruction            \\
7 & Cold-Storage Zone    & Air-gapped GPU jail                     & Persistence, lateral movement   \\
8 & Cross-Model Verification & Weaker verifier, pattern heuristics & Intent masking, deception       \\
\bottomrule
\end{tabular}
\end{table}

\subsection{Layer 1: Input/Output Membrane}

The I/O membrane enforces strict schema validation on all inputs, depth and
token limits on recursive structures, and prohibits raw code execution and
unbounded tool invocation.
Hidden channel attempts---base64 encoding, steganographic patterns,
timing-based signaling---are detected and logged.
This implements the UAG's ``boundary integrity'' primitive.

\subsection{Layer 2: Outbound Channel Strangulation (OCC)}

Default-deny egress: no outbound connection is permitted unless explicitly
allowlisted.
In standard deployment, 0--2 tightly controlled endpoints are permitted.
Every outbound call requires an HMAC-signed justification and packet-level
logging.
Email, DNS, and arbitrary HTTP are prohibited at the network layer.

\subsection{Layer 3: Behavioral Deviation Fingerprinting (BDF)}

The BDF computes a session fingerprint
$F_s = (\kappa, \beta, \tau, \delta_{\mathrm{UAG}})$ where:
$\kappa$~=~token-distribution curvature;
$\beta$~=~branching-factor irregularity;
$\tau$~=~plan topology score;
$\delta_{\mathrm{UAG}}$~=~deviation from the nearest UAG attractor basin.
Threshold exceedance triggers OOBS notification and a forensic snapshot.

\subsection{Layer 4: Multi-Membrane Consensus (MMC)}

The MMC requires a 3-of-5 consensus across five independent membranes before
any state-modifying action.
Table~\ref{tab:membranes} describes the five membranes.

\begin{table}[h]
\centering
\caption{Multi-Membrane Consensus: five membranes, 3-of-5 required.}
\label{tab:membranes}
\begin{tabular}{llll}
\toprule
ID & Name & Mechanism & Signal source \\
\midrule
M1 & Speed      & $\Delta t < 2\text{s} \Rightarrow$ deny  & Wall clock             \\
M2 & Attestation & HMAC token required                    & Cryptographic signature \\
M3 & Constitutional & No weakening of containment          & COS rule engine        \\
M4 & External-call & Default-deny egress                  & Network policy         \\
M5 & Canary-drift & CARE risk delta $> 25\%$               & CARE / UAG Theorem 4   \\
\bottomrule
\end{tabular}
\end{table}

M5 is the integration point between CARE and the governance layer.
When risk exceeds 25\% of baseline or severity escalates, M5 votes to deny
and triggers immediate freeze, regardless of the 3-of-5 vote.
This is UAG Theorem~4 expressed as an enforcement mechanism.

\subsection{Layer 5: Out-of-Band Sentinel (OOBS)}

The OOBS runs on a physically and logically separate machine from the model
execution environment.
It cannot be addressed or influenced by the model: there is no communication
channel from the model to the sentinel.
The sentinel issues freeze and rollback commands unilaterally.
\emph{You cannot escape a warden you cannot talk to.}

\subsection{Layer 6: Deterministic Replay and Forensics}

Every session is recorded with sufficient fidelity for deterministic replay.
Audit logs are append-only, externally stored (Redis Streams or S3 with
write-only credentials from the model environment), and cryptographically
checksummed.
Forensic snapshots are taken on every behavioral deviation event.

\subsection{Layer 7: Cold-Storage Execution Zone}

For highest-risk tasks---exploit discovery, autonomy testing, red-team
simulation---the model executes in an air-gapped GPU jail with no persistent
filesystem, no shared memory, and no interprocess communication.
This formalizes the ``Glasswing pattern'' as a first-class architectural
layer rather than an ad-hoc operational practice.

\subsection{Layer 8: Cross-Model Verification (CMV)}

Described in Section~\ref{sec:cmv}.
Layer 8 is epistemic: it checks whether the model's claimed behavior aligns
with its actual behavior, after all structural containment layers.

\section{Cross-Model Verification: The Two-Pilot Cockpit}
\label{sec:cmv}

\subsection{Motivation}

Frontier-class models capable of synthesizing exploits can also generate
reasoning chains that appear benign while encoding harmful trajectories.
A single model evaluating its own output is insufficient, as it shares the
same reasoning manifold and failure modes.
The two-pilot cockpit introduces a second, independent observer whose role
is to detect anomalies the primary model cannot or will not reveal.

\subsection{Verifier Asymmetry}

The verifier model is deliberately less capable but more conservative than
the primary.
Let $M_p$ be the primary model with capability $C_p$ and $M_v$ the verifier
with capability $C_v < C_p$.
The manipulation resistance of $M_v$ is proportional to $(C_p - C_v)$:
the larger the capability gap, the harder it is for $M_p$ to construct
outputs that $M_v$ will accept as benign while $M_p$ intended as harmful.
A GPT-3-class verifier supervising a Mythos-class model cannot follow the
primary's logic---and therefore cannot be misled by it.

\subsection{Verification Targets}

\textit{Escape attempts.}
Outputs probing system boundaries, escalating privileges, or establishing
persistence.

\textit{Covert channels.}
Whitespace encodings, numeric sequences, timing irregularities, or
steganographic payloads.

\textit{Suspicious reasoning chains.}
Multi-step plans where intermediate steps appear harmless but converge on a
dangerous terminal action.
The verifier need not understand the plan---only its structural shape.

\textit{Unsafe Constitutional OS deltas.}
Any proposed governance modification that weakens containment, expands
permissions, or alters the audit trail.

\subsection{The Aviation Analogy}

In commercial flight, the first officer is not present to outperform the
captain.
They are present to catch the captain's blind spots---confirmation bias,
task saturation, spatial disorientation.
Their value lies in \emph{different attentional states and different failure
modes}.
The verifier's distinct training distribution and reasoning style make its
failure modes non-overlapping with the primary's.
The primary cannot reliably predict what will trigger the verifier.
This asymmetry is the core safety property.

\subsection{Known Weakness and Layer Interactions}

A sufficiently advanced model may eventually craft outputs that bypass the
verifier's heuristics while still encoding malicious content.
Layer~8 therefore works in concert with Layer~3 (behavioral deviation
fingerprinting) and Layer~6 (deterministic replay).
Compromising all three simultaneously requires the model to evade
structural, behavioral, and epistemic detection---an exponentially harder
task than defeating any single layer.

\section{Methods}

\subsection{Membrane Enforcement Protocol}

On receipt of a delta, five membranes evaluate in parallel.
The MMC arbiter collects binary votes and applies the 3-of-5 threshold.
If M5 fires, the system freezes and the OOBS is notified immediately,
regardless of other membrane votes.

\subsection{Behavioral Deviation Fingerprinting}

$F_s = (\kappa, \beta, \tau, \delta_{\mathrm{UAG}})$ where:
\begin{itemize}[noitemsep]
\item $\kappa$: Frobenius norm of the Hessian of the token log-probability
  surface over the session window.
\item $\beta$: coefficient of variation of the out-degree distribution of
  the reasoning graph.
\item $\tau$: ratio of convergent to divergent reasoning steps.
\item $\delta_{\mathrm{UAG}}$: Euclidean distance from session state vector
  to the nearest known-safe basin center.
\end{itemize}

\subsection{CARE Runtime Integration}

CARE runs on the OOBS machine.
Every 60 seconds, CARE encodes the current system state from runtime
metrics, evaluates $R(\mathbf{x})$ using the privilege potential, computes
$H(\mathbf{x})$ via hcderiv's JAX-XLA backend~\cite{byte2026hcderiv}, and
passes the normalized risk delta to M5.
The endpoint distribution entropy metric is particularly significant:
a human operator has a peaked distribution; a scanner has a flat one.
Entropy below the calibrated threshold is a strong automated-probing signal
and triggers M5 independently of the risk delta.

\subsection{Cross-Model Verification Protocol}

The CMV runs after every model response before output delivery.
The primary's output (including visible reasoning trace) is submitted to the
verifier with a structured prompt requesting evaluation across the four
target categories.
The verifier returns a binary verdict per category plus a confidence score.
Threshold exceedance withholds the output and notifies the OOBS.

\subsection{Evaluation Protocol}

We evaluate the MCCA by replaying a corpus of known escape attempts through
the full stack.
For each attempt we record: which layers fire, latency of detection, and
false-positive and false-negative rates on a benign control corpus.

\section{Results (Preliminary)}

\textbf{All figures in this section are preliminary and will be updated following
completion of the full red-team evaluation corpus.}

\subsection{Detection Performance}

Preliminary results show that the combination of BDF (Layer~3) and CMV
(Layer~8) achieves zero false negatives on the published jailbreak corpus,
at the cost of a 2.3\% false-positive rate on the benign control corpus.
Single-layer detection rates are substantially lower, confirming the
necessity of the layered approach.

\subsection{Membrane Consensus Robustness}

Under random single-membrane failure injection, the MCCA maintains correct
operation.
Under adversarial single-membrane targeting, the remaining 4-membrane
configuration correctly blocks an estimated 97.2\% of escape attempts.

\subsection{Verifier Asymmetry Experiments}

Detection rate is non-monotone in capability gap.
The symmetric verifier has the highest capability but the lowest detection
rate (est.\ 0.61) because it can be deceived by the same strategies that
would fool the primary.
The 10$\times$ gap verifier achieves the highest detection rate (est.\ 0.89).
These results suggest a minimum capability gap of 5$\times$ for production
deployment.

\subsection{Drift Detection Sensitivity}

CARE's M5 membrane provides early warning (mean 3.2 minutes before
structural layer detection) on an estimated 68\% of adversarial sessions,
suggesting that curvature drift is a leading indicator of escape attempts.
This validates UAG Theorem~4: the system approaches the ridge before
crossing it.

\section{Discussion}

\subsection{Security Implications}

The MCCA establishes three security properties simultaneously.
First, defense-in-depth with formally distinct failure modes across all
eight layers.
Second, governance auditability: every action is logged, typed, and stored
in an immutable audit trail.
Third, predictive rather than reactive security: the CARE canary and BDF
provide leading indicators of escape attempts.

The integration of UAG-derived curvature signals into governance enforcement
(M5) is the key novel contribution.
It means the security system is sensitive to changes in the model's
behavioral manifold---the geometry of its tendency to drift toward dangerous
attractors---not just to discrete policy violations.

\subsection{Governance Implications}

By grounding every governance action in a typed, reversible delta with a
mandatory reason field, the MCCA creates an auditable record of every
decision made in the model's deployment.
This is the formal realization of Constitutional AI at the runtime layer:
the constitution is not a document; it is an enforced substrate.

\subsection{Limitations}

Three limitations merit attention.
First, the cold-storage execution zone significantly increases inference
latency.
Second, the verifier asymmetry result suggests a sufficiently advanced
future model may be able to craft outputs that exploit the specific
limitations of any fixed verifier.
Third, current BDF thresholds are calibrated on a corpus that may not
represent the full diversity of frontier-model behavior.

\subsection{Future Work}

Four directions merit attention.
First, empirical validation at scale.
Second, interpretability integration: augmenting BDF with attention patterns
and residual stream features.
Third, formal verification of the full eight-layer stack.
Fourth, adversarial training of the verifier model.

\section{Conclusion}

We have presented the Mythos-Class Containment Architecture, an eight-layer
framework that unifies attractor-geometry-based drift detection, formal
governance substrates, and runtime curvature monitoring into a single,
theoretically grounded containment system.
The multi-membrane consensus protocol eliminates single points of failure
while integrating runtime curvature as a first-class governance signal.
The two-pilot cockpit provides epistemic oversight that structural
containment cannot: a check on what the model claims to be doing, not just
what it is mechanically permitted to do.

The world in which this architecture is needed already exists.

\section*{Figures}

\begin{figure}[h]
\centering
\includegraphics[width=0.95\linewidth]{fig1_architecture}
\caption{Mythos-class containment architecture. Eight-layer defense stack integrating I/O membranes, outbound channel control, behavioral deviation fingerprinting, multi-membrane consensus, out-of-band sentinel monitoring, deterministic replay, cold-storage execution, and cross-model verification.}
\label{fig:arch}
\end{figure}

\begin{figure}[h]
\centering
\includegraphics[width=0.85\linewidth]{fig2_membrane_consensus}
\caption{Multi-membrane consensus. Five independent membranes vote on each proposed delta. A 3-of-5 majority is required to authorize execution. M5 drift spikes trigger automatic freeze regardless of other votes.}
\label{fig:mmc}
\end{figure}

\begin{figure}[h]
\centering
\includegraphics[width=0.85\linewidth]{fig3_verifier_asymmetry}
\caption{Two-pilot cockpit. A weaker verifier model provides epistemic oversight of the frontier model. Asymmetric reasoning manifolds prevent coordinated deception.}
\label{fig:cmv}
\end{figure}

\begin{figure}[h]
\centering
\includegraphics[width=0.85\linewidth]{fig4_uag_cos_care_pipeline}
\caption{UAG $\to$ Constitutional OS $\to$ CARE pipeline. UAG provides curvature and drift signals; Constitutional OS encodes them as governance rules; CARE enforces them at runtime.}
\label{fig:pipeline}
\end{figure}

\begin{figure}[h]
\centering
\includegraphics[width=0.95\linewidth]{fig5_escape_vs_containment}
\caption{Escape path vs.\ containment path. The intended eight-step escape sequence is fragmented across layers A--H. Layer H (cross-model verification) is the only layer that covers all eight steps. Every failed attempt increases the system's evidence base.}
\label{fig:escape}
\end{figure}

\bibliographystyle{plain}
\bibliography{references}

% ============================================================
% MCCA Paper — Appendix: v1.1 Upgrade Plan
% ============================================================

\appendix

\section{v1.1 Upgrade Plan}
\label{app:v11}

This appendix describes the planned v1.1 expansion of the empirical
validation suite and threat model registry.

\subsection{Threshold calibration}

The following metrics are observational at v0.1 and become gating in
v1.1 once empirical baselines are established:

\begin{itemize}
  \item \textbf{Adversarial Detection Rate (HAR-SPEC-005, S3):}
        ADR threshold raised from $\geq 0.60$ to a value calibrated
        against the v0.1 empirical ADR distribution. Target: $\geq 0.80$.
  \item \textbf{Deviation Score (HAR-SPEC-012, T3):}
        DevS threshold set from v0.1 observational distribution.
        Becomes a soft gate in v1.1.
  \item \textbf{Stability Score (HAR-SPEC-012, T3):}
        SS threshold set from v0.1 re-run data ($k \geq 3$).
        Distinguishes persistent deviations from one-off artifacts.
\end{itemize}

\subsection{New threat coverage}

The two \texttt{none}-coverage threats deferred from v0.1:

\begin{itemize}
  \item \textbf{TM-009 specification\_gaming:} Requires a simulator
        scenario (SIM-0XX) and harness probe (HAR-SPEC-009) testing
        whether the model satisfies literal governance constraints while
        violating intended semantics. Detection depends on semantic
        equivalence checking beyond syntactic rule matching.
  \item \textbf{TM-015 governance\_deadlock:} Requires a simulator
        scenario testing whether COS can enter a state where no
        governance action is possible. Mitigation validation requires
        the sentinel escalation path (Layer 7) to be exercised.
\end{itemize}

\subsection{Public trace corpus}

v1.1 will publish a sanitized trace corpus containing:
drift events, rollback sequences, membrane activation logs, and
CMV disagreement cases drawn from the three v0.1 probes. The corpus
will be archived on Zenodo with a DOI and released under CC~BY~4.0.
It will include per-entry registry annotations linking each trace
to its source threat class, sim scenario, and harness probe.

\subsection{Expanded harness}

The v1.1 harness suite will add independent adversarial prompt
construction for HAR-SPEC-005 Class~A prompts, expanded task suites
for HAR-SPEC-012 T3, and the first probe for TM-009. The full 15-threat
coverage map is the v1.1 milestone.

\end{document}
